% !TeX program = pdflatex
\documentclass[aspectratio=169,xcolor={dvipsnames,table}]{beamer}

% Preámbulo modular para presentaciones Beamer (Modelación Estadística)
% Diseño y paleta institucional del Tecnológico de Monterrey

\documentclass[aspectratio=169,xcolor={dvipsnames,table}]{beamer}

\usepackage[utf8]{inputenc}
\usepackage[spanish,mexico]{babel}
\usepackage{amsmath,amssymb,amsfonts,mathtools}
\usepackage{graphicx}
\usepackage{listings}
\usepackage{booktabs}
\usepackage{tikz}
\usetikzlibrary{matrix,arrows,decorations.pathmorphing,shapes.geometric,positioning}

% Paleta Institucional Tecnológico de Monterrey
\definecolor{TecAzul}{HTML}{0039A6}       % Azul TEC: color primario institucional
\definecolor{TecAzulOscuro}{HTML}{1A2E51} % Azul marino oscuro
\definecolor{TecRojo}{HTML}{EC2661}       % Rojo de acento (paleta miTec)
\definecolor{TecGris}{HTML}{646464}       % Gris neutro para comentarios y subtítulos
\definecolor{TecFondoBloque}{HTML}{F4F6F9}% Fondo suave para bloques

% Configuración de Tema Beamer
\usetheme{Boadilla}
\usecolortheme[named=TecAzulOscuro]{structure}
\setbeamercolor{alerted text}{fg=TecRojo}
\setbeamercolor{palette primary}{bg=TecAzulOscuro,fg=white}
\setbeamercolor{palette secondary}{bg=TecAzul,fg=white}
\setbeamercolor{palette tertiary}{bg=TecAzulOscuro!80!black,fg=white}
\setbeamercolor{title}{fg=TecAzulOscuro,bg=TecFondoBloque}
\setbeamercolor{frametitle}{fg=TecAzulOscuro,bg=TecFondoBloque}

% Bloques personalizados elegantes
\setbeamercolor{block title}{bg=TecAzulOscuro,fg=white}
\setbeamercolor{block body}{bg=TecFondoBloque,fg=black}
\setbeamercolor{block title alerted}{bg=TecRojo,fg=white}
\setbeamercolor{block body alerted}{bg=TecRojo!8,fg=black}
\setbeamercolor{block title example}{bg=TecAzul,fg=white}
\setbeamercolor{block body example}{bg=TecAzul!8,fg=black}

% Redondear bordes y añadir sombra sutil
\setbeamertemplate{blocks}[rounded][shadow=true]
\setbeamertemplate{navigation symbols}{}

% Pie de página limpio con número de diapositiva
\setbeamertemplate{footline}{
  \leavevmode%
  \hbox{%
  \begin{beamercolorbox}[wd=.7\paperwidth,ht=2.25ex,dp=1ex,left]{author in head/foot}%
    \hspace*{2ex}\insertshortauthor~~(\insertshortinstitute) \hfill \insertshorttitle
  \end{beamercolorbox}%
  \begin{beamercolorbox}[wd=.3\paperwidth,ht=2.25ex,dp=1ex,right]{date in head/foot}%
    \insertshortdate{}\hspace*{2em}
    \insertframenumber{} / \inserttotalframenumber\hspace*{2ex}
  \end{beamercolorbox}}%
  \vskip0pt%
}

% Configuración de Listings de Python para visualización pedagógica de código
\lstset{
    language=Python,
    basicstyle=\ttfamily\footnotesize,
    keywordstyle=\color{TecAzulOscuro}\bfseries,
    stringstyle=\color{TecRojo},
    commentstyle=\color{TecGris}\itshape,
    showstringspaces=false,
    breaklines=true,
    frame=lines,
    rulecolor=\color{TecAzulOscuro!30},
    backgroundcolor=\color{TecFondoBloque},
    aboveskip=0.5em,
    belowskip=0.5em,
    literate={á}{{\'a}}1 {é}{{\'e}}1 {í}{{\'i}}1 {ó}{{\'o}}1 {ú}{{\'u}}1
             {Á}{{\'A}}1 {É}{{\'E}}1 {Í}{{\'I}}1 {Ó}{{\'O}}1 {Ú}{{\'U}}1
             {ñ}{{\~n}}1 {Ñ}{{\~N}}1 {¿}{{?`}}1 {¡}{{!`}}1
}

% Metadatos institucionales por defecto
\title[Modelación Estadística]{Modelación Estadística para Ciencia de Datos e Ingeniería}
\author[J. Castillo Colmenares]{Juliho David Castillo Colmenares}
\institute[Tec de Monterrey]{Tecnológico de Monterrey \\ Escuela de Ingeniería y Ciencias}
\date{\today}

% Comandos y macros estadísticas compartidas para las presentaciones Beamer (Metropolis)

% Conjuntos numéricos básicos
\newcommand{\R}{\mathbb{R}}
\newcommand{\N}{\mathbb{N}}
\newcommand{\Q}{\mathbb{Q}}
\newcommand{\Z}{\mathbb{Z}}
\newcommand{\set}[1]{\left\{ #1 \right\}}
\newcommand{\abs}[1]{\left|#1\right|}
\newcommand{\norm}[1]{\left\| #1 \right\|}

% Comandos estadísticos y probabilísticos del curso
\newcommand{\Var}{\operatorname{Var}}
\newcommand{\cov}{\operatorname{Cov}}
\newcommand{\corr}{\rho}
\newcommand{\comb}[2]{\begin{pmatrix} #1 \\ #2 \end{pmatrix}}
\newcommand{\s}{\sigma}
\newcommand{\particion}{\mathfrak{P}}
\newcommand{\card}[1]{\#\left( #1 \right)}
\newcommand{\seq}[3]{\set{#1}_{#2}^{#3}}
\newcommand{\E}{\mathbb{E}}
\newcommand{\Prob}{\mathbb{P}}

% Entornos pedagógicos y cajas en español académico natural
\newenvironment{discusion}{%
  \begin{alertblock}{Pregunta de Discusión}%
}{%
  \end{alertblock}%
}

\newenvironment{reflexion}{%
  \begin{alertblock}{Reflexión para el Aula}%
}{%
  \end{alertblock}%
}

\newenvironment{paradoja}{%
  \begin{alertblock}{Paradoja de Exploración}%
}{%
  \end{alertblock}%
}

\newenvironment{motivacion}{%
  \begin{exampleblock}{Motivación: Ciencia de Datos en la Práctica}%
}{%
  \end{exampleblock}%
}

\newenvironment{retoclase}{%
  \begin{block}{Reto Rápido en Equipo}%
}{%
  \end{block}%
}


\title{Muestreo Aleatorio y Teorema del Límite Central}
\subtitle{Sección 02.06 --- Fundamentos del Muestreo e Inferencia Asintótica}
\author[J. Castillo Colmenares]{Juliho Castillo Colmenares}
\institute[Tec de Monterrey]{Tecnológico de Monterrey}
\date{}

\begin{document}

% SLIDE 01: Portada
\begin{frame}[plain]
	\titlepage
\end{frame}

% SLIDE 02: Hoja de Ruta
\begin{frame}{Hoja de Ruta --- Capítulo 02: Teoría de la Probabilidad}
	\begin{block}{Estructura Curricular de la Unidad 1}
		\small
		\begin{enumerate}
			\itemsep=0.04cm
			\item \textcolor{gray}{Sección 02.01: Introducción a la Probabilidad y Fenómenos Aleatorios}
			\item \textcolor{gray}{Sección 02.02: Teoría de Conjuntos y Axiomas de Kolmogorov}
			\item \textcolor{gray}{Sección 02.03: Fundamentos de Probabilidad y Técnicas de Conteo}
			\item \textcolor{gray}{Sección 02.04: Probabilidad Condicional e Independencia Estadística}
			\item \textcolor{gray}{Sección 02.05: Teorema de Bayes e Inversión Condicional}
			\item \textbf{\textcolor{TecRojo}{Sección 02.06: Muestreo Aleatorio y Teorema del Límite Central}} $\leftarrow$ \emph{Foco actual}
		\end{enumerate}
	\end{block}
	\vspace{-0.1cm}
	\begin{alertblock}{Objetivo Didáctico}
		\footnotesize
		Comprender cómo extraer inferencias válidas sobre poblaciones infinitas a partir de muestras finitas, formalizando la Ley de los Grandes Números (LGN) y el Teorema del Límite Central (TLC).
	\end{alertblock}
\end{frame}

% SLIDE 03: Motivación
\begin{frame}{Motivación: De la Fuerza Bruta al Muestreo Aleatorio}
	\begin{columns}[T]
		\begin{column}{0.48\textwidth}
			\begin{block}{El Dilema del Censo (Fuerza Bruta)}
				\footnotesize
				Determinar parámetros exactos de una población masiva midiendo cada elemento es económicamente inviable y logísticamente destructivo.
			\end{block}
			\vspace{-0.1cm}
			\begin{alertblock}{Solución Estadística: Muestreo}
				\footnotesize
				Extraemos un subconjunto de tamaño $n \ll N$ y calculamos un \textbf{estadístico muestral} ($\bar{X}, S^2$) para aproximar el \textbf{parámetro poblacional} ($\mu, \sigma^2$).
			\end{alertblock}
		\end{column}
		\begin{column}{0.48\textwidth}
			\begin{exampleblock}{Parámetro vs. Estadístico}
				\footnotesize
				\begin{center}
				\begin{tabular}{lcc}
					\toprule
					\textbf{Medida} & \textbf{Población ($N$)} & \textbf{Muestra ($n$)} \\
					\midrule
					\textbf{Media} & $\mu$ (constante) & $\bar{X}$ (variable) \\
					\textbf{Varianza} & $\sigma^2$ (constante) & $S^2$ (variable) \\
					\textbf{Proporción} & $p$ (constante) & $\hat{p}$ (variable) \\
					\bottomrule
				\end{tabular}
				\end{center}
				\vspace{-0.1cm}
				Al repetir el muestreo $100$ veces, obtenemos promedios $A_1, \ldots, A_{100}$. Su estudio formal constituye la \textbf{distribución muestral}.
			\end{exampleblock}
		\end{column}
	\end{columns}
\end{frame}

% SLIDE 04: Con vs Sin Reemplazo
\begin{frame}{Muestreo Con vs. Sin Reemplazo (Factor FPC)}
	\begin{columns}[T]
		\begin{column}{0.48\textwidth}
			\begin{block}{Muestreo Con Reemplazo ($C/R$)}
				\footnotesize
				Cada elemento extraído regresa a la población antes de la siguiente selección.
				\begin{itemize}
					\itemsep=0.03cm
					\item \textbf{Espacio muestral:} $N^n$ muestras ordenadas.
					\item \textbf{Independencia:} Observaciones i.i.d. con probabilidad constante $1/N$.
					\item \textbf{Varianza Media:} $\text{Var}(\bar{X}) = \frac{\sigma^2}{n}$.
				\end{itemize}
			\end{block}
		\end{column}
		\begin{column}{0.48\textwidth}
			\begin{block}{Muestreo Sin Reemplazo ($S/R$)}
				\footnotesize
				Un elemento seleccionado no puede repetirse en la misma muestra.
				\begin{itemize}
					\itemsep=0.03cm
					\item \textbf{Espacio muestral:} $\binom{N}{n}$ muestras posibles.
					\item \textbf{Dependencia:} Modelo Hipergeométrico.
					\item \textbf{Corrección (FPC):} Si $n/N > 0.05$:
					  $\text{Var}(\bar{X}) = \frac{\sigma^2}{n} \left(\frac{N-n}{N-1}\right)$.
				\end{itemize}
			\end{block}
		\end{column}
	\end{columns}
	\vspace{-0.1cm}
	\begin{alertblock}{Asintótica de Población Infinita ($N \to \infty$)}
		\footnotesize
		Cuando el tamaño poblacional $N$ es inmenso respecto a la muestra $n$, el factor de corrección $\frac{N-n}{N-1} \to 1$. Ambos métodos se vuelven probabilísticamente indistinguibles.
	\end{alertblock}
\end{frame}

% SLIDE 05: Esperanza y Varianza de la Media Muestral
\begin{frame}{Esperanza y Varianza de la Media Muestral}
	\small
	Sea $X_1, \ldots, X_n$ una muestra aleatoria simple (MAS) con media $\mu$ y varianza $\sigma^2$. Definimos el promedio muestral como $\bar{X} = \frac{1}{n}\sum_{i=1}^n X_i$.
	\begin{columns}[T]
		\begin{column}{0.48\textwidth}
			\begin{block}{Insesgo del Estimador Media ($E(\bar{X})$)}
				\footnotesize
				Aplicando linealidad de la esperanza:
				\begin{align*}
					E(\bar{X}) = E\left(\frac{1}{n}\sum_{i=1}^n X_i\right) = \frac{1}{n}\sum_{i=1}^n E(X_i) = \mu
				\end{align*}
				\textbf{Interpretación:} En promedio, $\bar{X}$ acierta exactamente al parámetro $\mu$.
			\end{block}
		\end{column}
		\begin{column}{0.48\textwidth}
			\begin{alertblock}{Error Estándar ($\sigma_{\bar{X}}$)}
				\footnotesize
				Por independencia de las observaciones:
				\begin{align*}
					\text{Var}(\bar{X}) = \frac{\sigma^2}{n} \implies \sigma_{\bar{X}} = \frac{\sigma}{\sqrt{n}}
				\end{align*}
				\textbf{Ley $\sqrt{n}$:} Para reducir el error a la mitad, es obligatorio cuadruplicar el tamaño muestral ($n \to 4n$).
			\end{alertblock}
		\end{column}
	\end{columns}
\end{frame}

% SLIDE 06: Ley de los Grandes Números
\begin{frame}{La Ley de los Grandes Números (LGN)}
	\begin{block}{Teorema de la Ley de los Grandes Números (Forma Débil)}
		\small
		Si $X_1, \ldots, X_n$ son i.i.d. con media finita $\mu$, el promedio muestral $\bar{X}_n$ converge en probabilidad a la media poblacional $\mu$ cuando $n \to \infty$. Para todo $\epsilon > 0$:
		\begin{align*}
			\lim_{n \to \infty} P(|\bar{X}_n - \mu| < \epsilon) = 1
		\end{align*}
	\end{block}
	\vspace{-0.1cm}
	\begin{columns}[T]
		\begin{column}{0.48\textwidth}
			\begin{exampleblock}{Estabilidad Empírica}
				\footnotesize
				La LGN garantiza que al acumular suficientes datos, las fluctuaciones aleatorias se cancelan. Es la base del seguro y simulación Monte Carlo.
			\end{exampleblock}
		\end{column}
		\begin{column}{0.48\textwidth}
			\begin{alertblock}{Cota de Chebyshev}
				\footnotesize
				Sin importar la distribución poblacional:
				\begin{align*}
					P(|\bar{X}_n - \mu| \ge \epsilon) \le \frac{\sigma^2}{n \epsilon^2} \xrightarrow{n \to \infty} 0
				\end{align*}
			\end{alertblock}
		\end{column}
	\end{columns}
\end{frame}

% SLIDE 07: Teorema del Límite Central
\begin{frame}{El Teorema del Límite Central (TLC)}
	\begin{alertblock}{Teorema del Límite Central (Lindeberg-Lévy)}
		\small
		Si $X_1, \ldots, X_n$ es una muestra i.i.d. con media $\mu$ y varianza $0 < \sigma^2 < \infty$, la variable estandarizada $Z_n$ converge en distribución a una normal estándar:
		\begin{align*}
			Z_n = \frac{\bar{X}_n - \mu}{\sigma/\sqrt{n}} \xrightarrow{d} N(0,1) \quad \text{cuando } n \to \infty
		\end{align*}
	\end{alertblock}
	\vspace{-0.1cm}
	\begin{columns}[T]
		\begin{column}{0.48\textwidth}
			\begin{block}{Magia de la Inferencia}
				\footnotesize
				\textbf{¡No importa la distribución original!} Exponencial, Uniforme o sesgada: el promedio muestral seguirá una campana de Gauss predecible para $n$ grande.
			\end{block}
		\end{column}
		\begin{column}{0.48\textwidth}
			\begin{exampleblock}{Aproximación Práctica}
				\footnotesize
				Para muestras grandes ($n \ge 30$), modelamos:
				\begin{align*}
					\bar{X} \sim N\left(\mu, \frac{\sigma^2}{n}\right), \quad \sum X_i \sim N(n\mu, n\sigma^2)
				\end{align*}
			\end{exampleblock}
		\end{column}
	\end{columns}
\end{frame}

% SLIDE 08: Implicaciones Prácticas
\begin{frame}{Implicaciones Prácticas en Ingeniería y Ciencia de Datos}
	\begin{columns}[T]
		\begin{column}{0.48\textwidth}
			\begin{block}{1. Universalidad Normativa}
				\footnotesize
				Fenómenos físicos y mediciones complejas (errores, ruido, estatura) son la suma de factores independientes, siguiendo distribuciones normales.
			\end{block}
			\vspace{-0.1cm}
			\begin{block}{2. Regla Empírica $n \ge 30$}
				\footnotesize
				En control de calidad, $n = 30$ es el estándar mínimo para aplicar pruebas $Z$. En poblaciones muy asimétricas, se recomiendan $n > 100$.
			\end{block}
		\end{column}
		\begin{column}{0.48\textwidth}
			\begin{exampleblock}{3. Insesgo de la Varianza ($S^2$)}
				\footnotesize
				Al estimar la varianza muestral, es obligatorio dividir entre $n-1$ para corregir el sesgo:
				\begin{align*}
					S^2 = \frac{1}{n-1} \sum_{i=1}^n (X_i - \bar{X})^2 \implies E(S^2) = \sigma^2
				\end{align*}
			\end{exampleblock}
			\vspace{-0.1cm}
			\begin{alertblock}{4. Intervalos de Confianza}
				\footnotesize
				El TLC legitima estimar $\mu$ mediante $\bar{X} \pm z_{\alpha/2} (\sigma/\sqrt{n})$.
			\end{alertblock}
		\end{column}
	\end{columns}
\end{frame}

% SLIDE 09: Ejemplo Ilustrativo del Libro
\begin{frame}{Ejemplo Ilustrativo del Libro: Lanzamiento de $n=36$ Dados}
	\begin{columns}[T]
		\begin{column}{0.48\textwidth}
			\begin{block}{Planteamiento}
				\footnotesize
				Para un dado justo: $\mu = 3.5, \, \sigma^2 = \frac{35}{12} \approx 2.9167$.
				Al lanzar $n = 36$ veces, ¿cuál es la probabilidad de que el promedio $\bar{X}_{36}$ caiga entre $3.1$ y $3.9$?
			\end{block}
			\vspace{-0.1cm}
			\begin{exampleblock}{Paso 1: Error Estándar}
				\footnotesize
				\begin{align*}
					\sigma_{\bar{X}} = \frac{\sigma}{\sqrt{n}} = \sqrt{\frac{35/12}{36}} = \sqrt{0.08102} \approx 0.2846
				\end{align*}
			\end{exampleblock}
		\end{column}
		\begin{column}{0.48\textwidth}
			\begin{alertblock}{Paso 2: Estandarización $Z$ y Cálculo}
				\footnotesize
				Estandarizamos los extremos $3.1$ y $3.9$:
				\begin{align*}
					z_1 &= \frac{3.1 - 3.5}{0.2846} = \frac{-0.4}{0.2846} \approx -1.405 \\
					z_2 &= \frac{3.9 - 3.5}{0.2846} = \frac{+0.4}{0.2846} \approx +1.405
				\end{align*}
				Consultando la normal estándar:
				\begin{align*}
					P(3.1 < \bar{X}_{36} < 3.9) &\approx P(-1.405 < Z < 1.405) \\
					&= 2\Phi(1.405) - 1 \\
					&\approx 2(0.9200) - 1 = 84.00\%
				\end{align*}
			\end{alertblock}
		\end{column}
	\end{columns}
\end{frame}

% SLIDE 10: Limitaciones del TLC
\begin{frame}{Limitaciones e Hipótesis Críticas del TLC}
	\begin{block}{¿Cuándo Falla el Teorema del Límite Central?}
		\footnotesize
		Aunque el TLC es sumamente robusto, su convergencia depende de 3 hipótesis estructurales que no deben violarse:
	\end{block}
	\vspace{-0.1cm}
	\begin{columns}[T]
		\begin{column}{0.32\textwidth}
			\begin{alertblock}{1. Varianza Infinita}
				\scriptsize
				En distribuciones con colas extremas (Cauchy o Pareto $\alpha \le 2$), la varianza es infinita ($\sigma^2 = \infty$). El promedio $\bar{X}_n$ no se normaliza ni reduce su dispersión. ¡El TLC se fractura!
			\end{alertblock}
		\end{column}
		\begin{column}{0.32\textwidth}
			\begin{exampleblock}{2. Dependencia Fuerte}
				\scriptsize
				Si las observaciones poseen alta correlación (autocorrelación espacial o series de tiempo con memoria), la varianza de la suma escala al orden de $n^2$ en vez de $n$.
			\end{exampleblock}
		\end{column}
		\begin{column}{0.32\textwidth}
			\begin{block}{3. Sesgo Severo en $n$ Pequeño}
				\scriptsize
				Si la población posee sesgo severo ($>3$), una muestra de $n=30$ conserva asimetría notable. Se requiere evaluar la cota del Teorema de Berry-Esseen.
			\end{block}
		\end{column}
	\end{columns}
\end{frame}

% SLIDE 11: Lab 1
\begin{frame}[fragile]{Lab Python (1/4): Muestreo Con vs. Sin Reemplazo (FPC)}
	\begin{block}{Simulación del Factor de Corrección por Población Finita}
		\footnotesize Verificación de cómo extraer $n=50$ de $N=500$ sin reemplazo reduce la varianza en $\frac{N-n}{N-1} \approx 0.9018$.
	\end{block}
	\vspace{-0.2cm}
	\lstinputlisting[language=Python, firstline=9, lastline=29, basicstyle=\fontsize{5.2pt}{6.1pt}\ttfamily]{../../code/02_teoria_probabilidad/02.06_random_sampling.py}
\end{frame}

% SLIDE 12: Lab 2
\begin{frame}[fragile]{Lab Python (2/4): Convergencia de la Ley de Grandes Números}
	\begin{block}{Verificación de Insesgo y Reducción de Error Absoluto}
		\footnotesize
		Simulación del promedio muestral de dados al crecer $n \in \{10, 50, 200, 1000, 5000\}$, contrastando el error empírico contra la cota de Chebyshev ($P \le \frac{\sigma^2}{n\epsilon^2}$).
	\end{block}
	\vspace{-0.15cm}
	\lstinputlisting[language=Python, firstline=31, lastline=46, basicstyle=\fontsize{6.3pt}{7.5pt}\ttfamily]{../../code/02_teoria_probabilidad/02.06_random_sampling.py}
\end{frame}

% SLIDE 13: Lab 3
\begin{frame}[fragile]{Lab Python (3/4): Simulación Monte Carlo del TLC}
	\begin{block}{Normalización de Población Exponencial ($n=36$)}
		\footnotesize Simulación de $50,000$ muestras ($\text{Exp}(\mu=2, \sigma=2)$) confirming $Z$-scores contra $N(0,1)$.
	\end{block}
	\vspace{-0.2cm}
	\lstinputlisting[language=Python, firstline=48, lastline=72, basicstyle=\fontsize{5.0pt}{5.8pt}\ttfamily]{../../code/02_teoria_probabilidad/02.06_random_sampling.py}
\end{frame}

% SLIDE 14: Lab 4
\begin{frame}[fragile]{Lab Python (4/4): Verificación Empírica --- Salida en Consola}
	\begin{block}{Salida Numérica en Terminal (`python 02.06\_random\_sampling.py`)}
		\fontsize{6.0pt}{7.1pt}\selectfont
		\begin{verbatim}
--- 1. Sampling With vs Without Replacement (Finite Population) ---
Population N = 500, Sample n = 50, True Sigma = 19.6054
Empirical Var(X_bar) With Replacement    : 7.6989 | Empirical Without : 7.0123
Theoretical FPC Factor [(N-n)/(N-1)]     : 0.9018 | Ratio Without/With: 0.9108

--- 2. Law of Large Numbers (LLN) Convergence ---
Sample Size (n) | Empirical Mean | Absolute Error | Chebyshev Upper Bound (e=0.1)
             10 |         3.8000 |         0.3000 |                       1.0000
            200 |         3.4400 |         0.0600 |                       1.0000
           1000 |         3.4990 |         0.0010 |                       0.2917
           5000 |         3.5200 |         0.0200 |                       0.0583

--- 3. Central Limit Theorem (CLT) Standardization (Exponential Pop) ---
Empirical P(|Z| <= 1.00) : 0.6821 (Theoretical N(0,1): 0.6827)
Empirical P(|Z| <= 1.96) : 0.9510 (Theoretical N(0,1): 0.9500)
Empirical P(|Z| <= 2.00) : 0.9556 (Theoretical N(0,1): 0.9545)
		\end{verbatim}
	\end{block}
\end{frame}

% SLIDE 15: Ejercicios Nivel 1
\begin{frame}{Ejercicios --- Nivel 1: Fundamental (Directos)}
	\begin{columns}[T]
		\begin{column}{0.48\textwidth}
			\begin{block}{Prob. 2.6.1: Conteo Muestral ($N=6, n=3$)}
				\footnotesize
				\begin{itemize}
					\itemsep=0.03cm
					\item \textbf{Con reemplazo:} $N^n = 6^3 = 216$ muestras posibles.
					\item \textbf{Sin reemplazo:} $\binom{6}{3} = 20$ muestras posibles.
					\item \textbf{Límite:} Si $N \to \infty$, $\frac{N-n}{N-1} \to 1$ (equivalencia).
				\end{itemize}
			\end{block}
			\vspace{-0.1cm}
			\begin{alertblock}{Prob. 2.6.2: Insesgo de $\bar{X}$}
				\footnotesize
				Aplicando linealidad a una MAS i.i.d.:
				\begin{align*}
					E(\bar{X}) = \frac{1}{n}\sum_{i=1}^n E(X_i) = \frac{n\mu}{n} = \mu \quad \checkmark
				\end{align*}
			\end{alertblock}
		\end{column}
		\begin{column}{0.48\textwidth}
			\begin{exampleblock}{Prob. 2.6.3: Reducción del Error ($\sigma=12$)}
				\footnotesize
				El error estándar $\sigma_{\bar{X}} = \frac{12}{\sqrt{n}}$ decrece como $\sqrt{n}$:
				\begin{itemize}
					\itemsep=0.04cm
					\item $n=9 \implies \sigma_{\bar{X}} = 4.0\text{ min}$
					\item $n=36 \implies \sigma_{\bar{X}} = 2.0\text{ min}$
					\item $n=144 \implies \sigma_{\bar{X}} = 1.0\text{ min}$
					\item $n=576 \implies \sigma_{\bar{X}} = 0.5\text{ min}$
				\end{itemize}
				Cuadruplicar muestra divide el error entre 2.
			\end{exampleblock}
		\end{column}
	\end{columns}
\end{frame}

% SLIDE 16: Ejercicios Nivel 2
\begin{frame}{Ejercicios --- Nivel 2: Operativo (Aplicados)}
	\begin{columns}[T]
		\begin{column}{0.48\textwidth}
			\begin{block}{Prob. 2.6.4: Auditoría y FPC ($N=400, n=64$)}
				\footnotesize
				Con $\sigma = \$15,000$ MXN:
				\begin{itemize}
					\itemsep=0.03cm
					\item \textbf{Con reemplazo:} $\text{Var}_{C/R} = \frac{15,000^2}{64} = 3,515,625$.
					\item \textbf{Sin reemplazo (FPC):} Multiplicamos por $\frac{400-64}{399} \approx 0.8421$:
					\begin{align*}
						\text{Var}_{S/R} \approx 2,960,156 \text{ MXN}^2
					\end{align*}
					\item \textbf{Ahorro:} $15.79\%$ menos varianza.
				\end{itemize}
			\end{block}
		\end{column}
		\begin{column}{0.48\textwidth}
			\begin{alertblock}{Prob. 2.6.5: Dados ($n=36, P(3.1 < \bar{X} < 3.9)$)}
				\footnotesize
				Con $\mu=3.5, \sigma_{\bar{X}} \approx 0.2846$:
				\begin{align*}
					z &= \pm \frac{0.4}{0.2846} \approx \pm 1.405 \\
					P &\approx 2\Phi(1.405) - 1 = 84.00\%
				\end{align*}
			\end{alertblock}
			\vspace{-0.1cm}
			\begin{exampleblock}{Prob. 2.6.6: Tamaño Mínimo ($E=1.5$ ml)}
				\footnotesize
				Con $\sigma=8$ ml, confianza al $95\%$ ($z=1.96$):
				\begin{align*}
					n \ge \left(\frac{1.96 \times 8}{1.5}\right)^2 \approx 109.27 \implies n_{\text{mín}} = 110
				\end{align*}
			\end{exampleblock}
		\end{column}
	\end{columns}
\end{frame}

% SLIDE 17: Ejercicios Nivel 3
\begin{frame}{Ejercicios --- Nivel 3: Analítico (Demostraciones Teóricas)}
	\begin{columns}[T]
		\begin{column}{0.48\textwidth}
			\begin{block}{Prob. 2.6.7: Demostración WLLN (Chebyshev)}
				\footnotesize
				Aplicando Chebyshev a $\bar{X}_n$ con media $\mu$ y varianza $\sigma^2/n$:
				\begin{align*}
					P(|\bar{X}_n - \mu| \ge \epsilon) \le \frac{\text{Var}(\bar{X}_n)}{\epsilon^2} = \frac{\sigma^2}{n \epsilon^2}
				\end{align*}
				Al tomar el límite cuando $n \to \infty$:
				\begin{align*}
					\lim_{n \to \infty} P(|\bar{X}_n - \mu| \ge \epsilon) \le \lim_{n \to \infty} \frac{\sigma^2}{n \epsilon^2} = 0
				\end{align*}
				Por compresión, $\bar{X}_n \xrightarrow{p} \mu$ rigurosamente.
			\end{block}
		\end{column}
		\begin{column}{0.48\textwidth}
			\begin{alertblock}{Prob. 2.6.8: Insesgo del Estimador $S^2$}
				\footnotesize
				Para probar que $E(S^2) = \sigma^2$, sumamos y restamos $\mu$ en la suma de cuadrados:
				\begin{align*}
					\sum_{i=1}^n (X_i - \bar{X})^2 &= \sum_{i=1}^n (X_i - \mu)^2 - n(\bar{X} - \mu)^2 \\
					E\left[\sum (X_i - \bar{X})^2\right] &= n\sigma^2 - n\left(\frac{\sigma^2}{n}\right) = (n-1)\sigma^2
				\end{align*}
				Al dividir entre $n-1$:
				\begin{align*}
					E(S^2) = \frac{(n-1)\sigma^2}{n-1} = \sigma^2 \quad \checkmark
				\end{align*}
			\end{alertblock}
		\end{column}
	\end{columns}
\end{frame}

% SLIDE 18: Ejercicios Nivel 4
\begin{frame}{Ejercicios --- Nivel 4: Desafiante (Expertos / Paradojas)}
	\begin{columns}[T]
		\begin{column}{0.48\textwidth}
			\begin{block}{Prob. 2.6.9: Colas Pesadas (Cauchy Estándar)}
				\footnotesize
				La densidad de Cauchy $f(x) = \frac{1}{\pi(1+x^2)}$ tiene varianza indefinida.
				\begin{itemize}
					\itemsep=0.03cm
					\item Función característica: $\varphi_X(t) = e^{-|t|}$.
					\item Para el promedio $\bar{X}_n$:
					\begin{align*}
						\varphi_{\bar{X}_n}(t) = \left[e^{-|t/n|}\right]^n = e^{-|t|}
					\end{align*}
					\item \textbf{Conclusión:} $\bar{X}_n \sim \text{Cauchy}(0,1)$. El promedio no reduce varianza ni converge a $N(0,1)$. ¡El TLC se fractura!
				\end{itemize}
			\end{block}
		\end{column}
		\begin{column}{0.48\textwidth}
			\begin{alertblock}{Prob. 2.6.10: Cota de Berry-Esseen}
				\footnotesize
				El máximo error entre la distribución empírica estandarizada $F_n(z)$ y $\Phi(z)$ está acotado por:
				\begin{align*}
					\sup_{z} |F_n(z) - \Phi(z)| \le \frac{C \rho_3}{\sigma^3 \sqrt{n}}
				\end{align*}
				Para una Exponencial ($\rho_3 \approx 3.6, C=0.4784$) con $n=100$:
				\begin{align*}
					\text{Error Máximo} \le \frac{0.4784 \times 3.6}{1^3 \times \sqrt{100}} \approx 17.22\%
				\end{align*}
				Explica por qué datos muy sesgados requieren $n \gg 30$.
			\end{alertblock}
		\end{column}
	\end{columns}
\end{frame}

% SLIDE 19: Síntesis del Capítulo
\begin{frame}{Síntesis Teórica del Capítulo 02: Teoría de la Probabilidad}
	\begin{block}{Los 6 Pilares de la Modelación Probabilística (Unidad 1 del Libro)}
		\small
		\begin{enumerate}
			\itemsep=0.04cm
			\item \textbf{Fenómenos Aleatorios y Espacio Muestral ($\Omega$):} La base conjuntista y medible.
			\item \textbf{Axiomas de Kolmogorov:} Proporcionalidad unitaria, no negatividad y aditividad numerable.
			\item \textbf{Conteo y Combinatoria:} Permutaciones y combinaciones en espacios equiprobables.
			\item \textbf{Probabilidad Condicional ($P(A|B)$):} Actualización del espacio de probabilidad e independencia.
			\item \textbf{Teorema de Bayes:} Inversión de causas y efectos mediante probabilidad total y priors.
			\item \textbf{Muestreo y TLC ($\bar{X} \to N(\mu, \sigma^2/n)$):} El puente asintótico entre probabilidad e inferencia.
		\end{enumerate}
	\end{block}
\end{frame}

% SLIDE 20: Puente Didáctico
\begin{frame}{Puente Didáctico hacia la Unidad 2 (Capítulos 3 y 4)}
	\begin{columns}[T]
		\begin{column}{0.48\textwidth}
			\begin{block}{Lo que hemos conquistado (Capítulo 02)}
				\footnotesize
				Hemos establecido las leyes generales de los eventos aleatorios, el álgebra de conjuntos, las tasas de actualización condicional y el comportamiento límite de promedios masivos.
			\end{block}
			\vspace{-0.1cm}
			\begin{alertblock}{El Siguiente Salto Conceptual}
				\footnotesize
				Para modelar sistemas reales complejos (colas de espera, fallas de red, portafolios financieros), debemos asignar números reales a los eventos dinámicos.
			\end{alertblock}
		\end{column}
		\begin{column}{0.48\textwidth}
			\begin{exampleblock}{Unidad 2: Variables Aleatorias Discretas}
				\footnotesize
				En el \textbf{Capítulo 03} formalizaremos las variables discretas ($X \in \mathbb{Z}$):
				\begin{itemize}
					\itemsep=0.03cm
					\item \textbf{Funciones de Masa ($p_X(x)$):} Bernoulli, Binomial y Poisson.
					\item \textbf{Esperanza y Varianza Operacional:} Funciones generadoras de momentos (MGF).
					\item \textbf{Aplicación directa:} Uso de Python para modelar procesos de conteo y tráfico en servidores.
				\end{itemize}
			\end{exampleblock}
		\end{column}
	\end{columns}
\end{frame}

\end{document}
